\begin{frame}{First Order Logic}
\begin{center}
\vspace{1.5cm}
{\Huge\textbf{First Order Logic}}
\end{center}
\end{frame}

% --------------------------------------------------
% Question 1
% --------------------------------------------------

\begin{frame}{Question 1}
A university maintains a database of students enrolled in a programming
course.

Let:

$$
P(x) = \text{$x$ has completed the programming course}
$$

Express the following requirement using first order logic:

\begin{center}
\textit{Every student enrolled in the course has completed the programming
course.}
\end{center}
\end{frame}

\begin{frame}{Question 1 -- Answer}
The statement applies to every student in the course.

Therefore, we use the universal quantifier:

$$
\forall x
$$

The predicate is:

$$
P(x) = \text{$x$ has completed the programming course}
$$

\end{frame}

\begin{frame}{Question 1 -- Answer}
The required statement is:

$$
\boxed{\forall x\,P(x)}
$$

This means that for every student $x$ in the relevant domain,
$x$ has completed the programming course.
\end{frame}

% --------------------------------------------------
% Question 2
% --------------------------------------------------

\begin{frame}{Question 2}
A cloud computing system contains several servers. The system requirement
states:

\begin{center}
\textit{At least one server must be operational before a service can be
provided.}
\end{center}

Let:

$$
P(x) = \text{$x$ is an operational server}
$$

Express this requirement using first order logic.
\end{frame}

\begin{frame}{Question 2 -- Answer}
The phrase \textit{at least one} indicates that we need an existential
quantifier.

The existential quantifier is:

$$
\exists x
$$

The predicate is:

$$
P(x) = \text{$x$ is an operational server}
$$

\end{frame}

\begin{frame}{Question 2 -- Answer}
Therefore, the requirement can be written as:

$$
\boxed{\exists x\,P(x)}
$$

This states that there exists at least one server that is operational.
\end{frame}

% --------------------------------------------------
% Question 3
% --------------------------------------------------

\begin{frame}{Question 3}
A company's authentication system records user login attempts.

Let:

$$
P(x) = \text{$x$ is a successful login}
$$

$$
Q(x) = \text{$x$ is associated with an authorized user}
$$

The security requirement states:

\begin{center}
\textit{Every successful login must be associated with an authorized user.}
\end{center}

Express the requirement using first order logic.
\end{frame}

\begin{frame}{Question 3 -- Answer}
The statement applies to every login attempt.

Therefore, we use:

$$
\forall x
$$

The condition that must hold for every successful login is:

$$
P(x) \rightarrow Q(x)
$$

\end{frame}

\begin{frame}{Question 3 -- Answer}
Combining the universal quantifier with the implication gives:

$$
\boxed{\forall x\,[P(x)\rightarrow Q(x)]}
$$

This means that whenever a login is successful, it must be associated
with an authorized user.
\end{frame}

% --------------------------------------------------
% Question 4
% --------------------------------------------------

\begin{frame}{Question 4}
An artificial intelligence system processes a collection of training
records.

Let:

$$
P(x) = \text{$x$ contains missing values}
$$

$$
Q(x) = \text{$x$ is removed from the training dataset}
$$

The preprocessing requirement states:

\begin{center}
\textit{Every record containing missing values is removed from the
training dataset.}
\end{center}

Express the requirement using first order logic.
\end{frame}

\begin{frame}{Question 4 -- Answer}
The statement applies to every training record.

Therefore, we use:

$$
\forall x
$$

If a record contains missing values, it must be removed.

This gives the implication:

$$
P(x)\rightarrow Q(x)
$$

\end{frame}

\begin{frame}{Question 4 -- Answer}
Therefore, the complete first order logic expression is:

$$
\boxed{\forall x\,[P(x)\rightarrow Q(x)]}
$$

The expression states that every record satisfying $P(x)$ must also
satisfy $Q(x)$.
\end{frame}

% --------------------------------------------------
% Question 5
% --------------------------------------------------

\begin{frame}{Question 5}
A software testing system contains a collection of test cases.

Let:

$$
P(x) = \text{$x$ exposes a software defect}
$$

The testing report claims:

\begin{center}
\textit{No test case exposes a software defect.}
\end{center}

Express this statement using first order logic and give an equivalent
form using an existential quantifier.
\end{frame}

\begin{frame}{Question 5 -- Answer}
The statement says that every test case does not expose a defect.

Therefore, we can use the universal quantifier:

$$
\forall x
$$

The statement becomes:

$$
\boxed{\forall x\,\neg P(x)}
$$

\end{frame}

\begin{frame}{Question 5 -- Answer}
The statement can also be expressed by saying that there does not exist
a test case that exposes a defect.

Using the existential quantifier:

$$
\boxed{\neg\exists x\,P(x)}
$$

\end{frame}

\begin{frame}{Question 5 -- Answer}
Therefore, the two equivalent forms are:

$$
\forall x\,\neg P(x)
$$

and

$$
\neg\exists x\,P(x)
$$

Hence:

$$
\boxed{\forall x\,\neg P(x)
\equiv
\neg\exists x\,P(x)}
$$

Both statements mean that no test case exposes a software defect.
\end{frame}
